\chapter{Краткая памятка}

\section{Системы счисления и переводы}
В позиционной системе значение цифры зависит от её места. В записи целого неотрицательного числа
\[
    (a_{m-1}\ldots a_1a_0)_p=\sum_{i=0}^{m-1}a_i p^i
\]
основание $p\ge2$, а каждая цифра удовлетворяет $0\le a_i<p$. В непозиционной системе нет такого набора позиционных весов: например, в унарной записи значение определяется количеством одинаковых знаков.

В шестнадцатеричной системе A, B, C, D, E, F обозначают 10, 11, 12, 13, 14, 15. Поэтому $19_8$ --- недопустимая запись: цифры 9 в восьмеричной системе нет.

\subsection*{Из десятичной системы в основание $p$}
Делите число на $p$ до нулевого частного, выписывая остатки. Ответ состоит из остатков \textbf{в обратном порядке}. Ноль переводится в ноль.

\subsection*{Из основания $p$ в десятичную систему}
Умножьте каждую цифру на вес её разряда и сложите. Веса разрядов справа налево: $1, p, p^2, \ldots$.

\subsection*{Быстрый перевод}
Основания $4=2^2$, $8=2^3$ и $16=2^4$ позволяют заменять одну цифру группой из 2, 3 или 4 битов. Целую двоичную запись группируйте \textbf{справа налево}; неполную старшую группу дополняйте нулями слева. При обратном переводе сохраняйте полную ширину каждой внутренней группы.

\begin{center}
\begin{tabular}{c|c@{\qquad}c|c}
Hex & Биты & Hex & Биты\\\hline
0 & \texttt{0000} & 8 & \texttt{1000}\\
1 & \texttt{0001} & 9 & \texttt{1001}\\
2 & \texttt{0010} & A & \texttt{1010}\\
3 & \texttt{0011} & B & \texttt{1011}\\
4 & \texttt{0100} & C & \texttt{1100}\\
5 & \texttt{0101} & D & \texttt{1101}\\
6 & \texttt{0110} & E & \texttt{1110}\\
7 & \texttt{0111} & F & \texttt{1111}
\end{tabular}
\end{center}

\begin{center}
\begin{tabular}{c|rrrrrrr}
$k$ & 0 & 1 & 2 & 3 & 4 & 5 & 6\\\hline
$2^k$ & 1 & 2 & 4 & 8 & 16 & 32 & 64\\[2pt]
$k$ & 7 & 8 & 9 & 10 & 11 & 12 & 16\\\hline
$2^k$ & 128 & 256 & 512 & 1024 & 2048 & 4096 & 65536
\end{tabular}
\end{center}

Дробные части переводят последовательным умножением на основание, читая полученные целые части в прямом порядке. В этом практикуме основные вычислительные задания посвящены целым числам.

\section{Арифметика в основании \texorpdfstring{$p$}{p}}
\begin{description}
    \item[Сложение.] Суммируйте цифры справа налево с учётом входящего переноса. Цифра результата --- остаток суммы от деления на $p$, новый перенос --- целая часть от деления на $p$.
    \item[Вычитание.] Если текущей цифры не хватает, займите единицу из старшего разряда: она даёт $p$ единиц текущего разряда. При необходимости заём проходит через несколько нулей.
    \item[Умножение.] Умножайте на каждую цифру множителя, переносите полные группы по $p$ и сдвигайте частичные произведения на разряд влево.
    \item[Деление.] Подбирайте цифры частного слева направо, вычитая кратные делителю. Для положительного делителя $B$ проверка имеет вид $A=BQ+R$, где $0\le R<B$.
\end{description}
Для независимой проверки переведите операнды и результат в десятичную систему. Если в условии требуется действие в заданном основании, покажите разрядные действия в этом основании.

\section{Разрядность и диапазоны}
$n$ битов дают $2^n$ различных комбинаций. Их числовые значения зависят от выбранного кода.

\begin{center}
\small
\begin{tabular}{l|c|c}
Представление & Диапазон & Различных значений\\\hline
Беззнаковое & $0\ldots 2^n-1$ & $2^n$\\
Прямой код & $-(2^{n-1}-1)\ldots 2^{n-1}-1$ & $2^n-1$\\
Обратный код & $-(2^{n-1}-1)\ldots 2^{n-1}-1$ & $2^n-1$\\
Дополнительный & $-2^{n-1}\ldots 2^{n-1}-1$ & $2^n$\\
Со смещением $b$ & $-b\ldots 2^n-1-b$ & $2^n$
\end{tabular}
\end{center}

У прямого и обратного кодов два представления нуля, но одно числовое значение 0. Поэтому число значений на единицу меньше числа битовых комбинаций.

Чтобы выбрать минимальную разрядность, проверьте \textbf{обе границы} нужного диапазона. Для беззнакового максимума $M$ требуется $M\le2^n-1$. Для диапазона $[L,U]$ в дополнительном коде нужны одновременно
\[
    -2^{n-1}\le L,\qquad U\le2^{n-1}-1.
\]
Обоснуйте минимальность: при $n-1$ битах хотя бы одна граница уже не помещается.

\section{Четыре способа кодирования}
Обозначим через $u$ беззнаковое значение битового слова, через $x$ --- кодируемое число.

\subsection*{Прямой код (sign-magnitude)}
Старший бит задаёт знак, остальные $n-1$ битов --- модуль. Для отрицательного числа код равен $2^{n-1}+|x|$. При чтении отделите знак от модуля. В байте нули имеют коды \texttt{00000000} и \texttt{10000000}.

\subsection*{Обратный код (ones' complement)}
Для $x<0$ инвертируйте все $n$ битов положительного числа $|x|$. Чтобы прочитать отрицательный код, инвертируйте слово и поставьте минус перед полученным модулем. Два нуля: все нули и все единицы.

\subsection*{Дополнительный код (two's complement)}
Для $x<0$ код равен $2^n+x$: можно инвертировать $n$-битную запись модуля и прибавить единицу, сохранив $n$ младших битов. Для $x\ge0$ код совпадает с обычной двоичной записью.

Чтение кода:
\[
x=\begin{cases}
u,&u<2^{n-1},\\
u-2^n,&u\ge2^{n-1}.
\end{cases}
\]
Эквивалентно, старший бит имеет вес $-2^{n-1}$, остальные --- обычные положительные веса. \textbf{Нельзя} просто удалить старшую единицу и назвать оставшиеся биты модулем.

Минимальное число $-2^{n-1}$ имеет код \texttt{100\ldots0}. Положительное $2^{n-1}$ уже не помещается в тот же знаковый формат. Поэтому смена знака минимального числа в фиксированной сетке вызывает знаковое переполнение.

\subsection*{Код со смещением (biased representation)}
Храним $u=x+b$, читаем $x=u-b$, где $b$ --- известное смещение. Старший бит здесь \textbf{не является обычным знаковым битом}. При $b=128$ и ширине 8 бит ноль имеет код \texttt{10000000}.

При сложении двух смещённых кодов смещение учтено дважды. Код суммы вычисляется так:
\[
u_{A+B}=u_A+u_B-b.
\]
Сначала вычислите точную сумму значений и проверьте диапазон. Коррекция кода не делает непредставимую сумму представимой.

\section{Ограниченная разрядность и флаги}
Для $n$-битного слова младшие биты суммы или разности вычисляются по модулю $2^n$. Например, в 8 битах остаётся остаток в диапазоне $0\ldots255$.

Одинаковые биты результата можно читать беззнаково и знаково. Точная математическая сумма может отличаться от прочитанного значения, если произошёл выход за соответствующий диапазон.

\begin{description}
    \item[CF при сложении.] Равен 1, если сумма беззнаковых значений операндов не меньше $2^n$; иначе 0.
    \item[CF при вычитании x86.] Равен 1, если из беззнакового $u_A$ вычитают большее $u_B$, то есть нужен заём; иначе 0.
    \item[OF.] Равен 1, если точный результат операции над знаковыми значениями не помещается в диапазон дополнительного кода.
    \item[ZF.] Равен 1, если все сохранённые биты результата нулевые.
    \item[SF.] Копирует старший бит сохранённого результата; при переполнении это не обязательно знак точного математического ответа.
\end{description}

Быстрые правила OF:
\begin{itemize}
    \item При сложении: одинаковые знаки операндов, но другой знак результата.
    \item При вычитании: разные знаки операндов, а знак результата отличается от знака уменьшаемого.
\end{itemize}
CF и OF описывают разные интерпретации. Один может быть равен 1, а другой 0; возможны и две единицы, и два нуля.

Вычитание можно вычислить как $A+(\mathop{\mathrm{NOT}}B+1)$, сохранив младшие $n$ битов. Но \textbf{перенос этого вспомогательного сложения не подставляют вместо CF инструкции SUB}. Для SUB проверяют $u_A<u_B$. В частности, одинаковые биты результата не означают одинаковые флаги у разных инструкций.

Эта модель объясняет работу процессора. В языке C беззнаковая арифметика подчиняется модулю соответствующего типа, а знаковое переполнение нельзя использовать как гарантированный способ получить «обёрнутый» результат: оно имеет неопределённое поведение.

\section{Побитовые операции и маски}
\begin{center}
\begin{tabular}{cc|ccc}
$a$ & $b$ & AND & OR & XOR\\\hline
0 & 0 & 0 & 0 & 0\\
0 & 1 & 0 & 1 & 1\\
1 & 0 & 0 & 1 & 1\\
1 & 1 & 1 & 1 & 0
\end{tabular}
\end{center}
NOT меняет каждый бит на противоположный. Для $n$-битного слова:
\[
\mathop{\mathrm{NOT}}x=(2^n-1)-x.
\]
Результат инверсии зависит от ширины слова!

В записи, похожей на C: AND --- \texttt{\&}, OR --- \texttt{|}, XOR --- \texttt{\^{}}, NOT --- \texttt{\textasciitilde}. Пусть $m=2^k$ --- маска с единицей только в бите $k$. Тогда:
\begin{itemize}
    \item проверить бит: \texttt{(x \& m) != 0}; результат \texttt{x \& m} сам по себе не обязан быть равен 1;
    \item установить бит: \texttt{x | m};
    \item сбросить бит: \texttt{x \& NOT(m)}, где NOT ограничен шириной слова;
    \item инвертировать бит: \texttt{x XOR m}.
\end{itemize}
Чтобы выделить $w$ битов начиная с номера $k$, выполните логический сдвиг вправо на $k$, затем AND с маской $2^w-1$. Чтобы заменить поле, сначала очистите его старое содержимое, затем установите новое через OR.

\section{Сдвиги и изменение разрядности}
Для всех сдвигов здесь задано $0<k<n$.
\begin{description}
    \item[Логический влево.] Вышедшие слева биты теряются, справа добавляются нули. Для unsigned результат равен $x2^k\bmod2^n$. Точное произведение сохраняется, если оно помещается в $n$ битов.
    \item[Логический вправо.] Слева добавляются нули; результат равен $\lfloor u/2^k\rfloor$.
    \item[Арифметический вправо.] Слева добавляются копии знакового бита. Для дополнительного кода результат равен $\lfloor x/2^k\rfloor$, то есть округляется вниз, в том числе для отрицательного $x$.
    \item[Циклический.] Вытесненные биты возвращаются с другого края.
\end{description}

Например, округление вниз и к нулю для отрицательного нецелого частного различаются. Деление знаковых целых в C округляет к нулю. Поведение сдвига отрицательного signed в C зависит от применяемых правил стандарта и реализации; когда нужен именно арифметический сдвиг, это следует задавать явно. В NASM/x86 различают \texttt{SHR}, \texttt{SAR}, \texttt{ROL}, \texttt{ROR}.

\subsection*{Расширение}
При \textbf{нулевом расширении} unsigned добавьте слева нули. При \textbf{знаковом расширении} дополнительного кода повторяйте старший бит. Значение числа сохраняется, но выбранный способ зависит от интерпретации исходного слова.

\subsection*{Сужение}
Отбрасывание старших битов не всегда сохраняет значение. Проверяйте, помещается ли исходное число в диапазон нового формата. Сначала укажите оставшийся код, затем прочитайте его в требуемой интерпретации.
